\chapter{Implémentation --- Frontend Angular}

\section*{Introduction}
Ce chapitre présente l'implémentation de l'interface utilisateur de la plateforme PMS, développée
avec \textbf{Angular~21}. Après avoir exposé l'architecture retenue (composants \emph{standalone},
API \emph{Signals}, structure par fonctionnalités), nous détaillons les mécanismes transversaux~:
authentification JWT côté client, application du contrôle d'accès dynamique (RBAC) à l'interface,
communication avec l'API REST, puis les écrans métier et le reflet du périmètre de données
(\emph{scope}) dans l'affichage.

\section{Architecture de l'application frontend}

L'application est une \emph{Single Page Application} (SPA) bâtie sur Angular~21 en mode
\textbf{standalone} (sans \texttt{NgModule}), avec l'API \textbf{Signals} pour la gestion d'état
réactive, \textbf{Bootstrap~5} et \texttt{bootstrap-icons} pour la présentation, et une interface
entièrement en \textbf{français} (ADR-006).

La structure de dossiers reproduit le découpage par fonctionnalités du backend (ADR-016), ce qui
facilite la navigation et la traçabilité entre une fonctionnalité métier et son écran~:

\begin{verbatim}
src/app/
  core/
    models/        interfaces typees (Project, Mission, ...)
    services/      acces API (ProjectService, AuthService, ...)
    guards/        authGuard, permissionGuard
    interceptors/  ajout du jeton JWT, rafraichissement
  features/
    auth/          login, changement de mot de passe
    dashboard/     tableau de bord
    projects/      liste, detail (6 onglets), formulaire
    workload/  kpi/  billing/  missions/  governance/
    resources/     tarifs journaliers (TCC)
    admin/         gestion des utilisateurs
  layout/shell/    cadre applicatif + menu lateral dynamique
\end{verbatim}

\noindent Chaque écran est un composant autonome chargé en \emph{lazy loading} via
\texttt{loadComponent}, ce qui réduit la taille du \emph{bundle} initial et accélère le premier
affichage (NFR-001).

\section{Authentification et session côté client}

À la connexion, le composant de \emph{login} appelle \texttt{POST /api/auth/login}~; la réponse
enrichie contient le jeton d'accès, le jeton de rafraîchissement et le \textbf{contexte utilisateur}
(courriel, nom complet, rôle et \textbf{ensemble des permissions}). Ce contexte est conservé par
\texttt{AuthService} et constitue la source unique des décisions d'affichage.

Un \textbf{intercepteur HTTP} ajoute automatiquement l'en-tête \texttt{Authorization: Bearer~<token>}
à chaque requête sortante vers l'API. En cas de réponse \texttt{401}, il déclenche un appel à
\texttt{POST /api/auth/refresh} puis rejoue la requête~; si le rafraîchissement échoue, l'utilisateur
est redirigé vers l'écran de connexion.

Le drapeau \texttt{firstLogin} impose le changement du mot de passe temporaire~: tant qu'il vaut
\texttt{true}, un garde de route force la redirection vers \texttt{/change-password} et bloque l'accès
au reste de l'application (ADR-010).

\section{Application du RBAC dynamique à l'interface (ADR-001)}

Le principe directeur du projet --- l'autorisation pilotée par \textbf{permission} et jamais par nom
de rôle --- est appliqué de bout en bout, y compris dans l'interface. Trois niveaux coexistent~:

\begin{itemize}
  \item \textbf{Protection des routes}~: chaque route sensible est gardée par \texttt{authGuard}
        (utilisateur authentifié) et \texttt{permissionGuard}, qui lit la permission requise depuis
        les données de la route et la confronte au contexte utilisateur.
  \item \textbf{Menu latéral dynamique}~: les entrées du menu sont générées à partir de l'ensemble
        des permissions reçu à la connexion. Un utilisateur ne voit que les modules auxquels il a
        droit (le développeur ne voit ni \emph{Facturation} ni \emph{Administration}).
  \item \textbf{Masquage des actions}~: à l'intérieur d'un écran, les boutons d'écriture
        (créer, modifier, supprimer, assigner) sont conditionnés par \texttt{auth.hasPermission(\dots)}.
\end{itemize}

\begin{verbatim}
// Extrait — garde de permission appliqué à une route
{
  path: 'billing',
  canActivate: [permissionGuard],
  data: { permission: 'VIEW_BILLING' },
  loadComponent: () => import('./features/billing/billing.component')
                         .then(m => m.BillingComponent)
}
\end{verbatim}

\noindent Cette protection côté client est un \textbf{confort d'usage}~: elle évite de présenter des
actions interdites. L'\textbf{autorité} reste le backend, qui revérifie systématiquement chaque appel
par \texttt{@PreAuthorize}. Un utilisateur qui contournerait l'interface se heurterait à une réponse
\texttt{403}. Cette double application illustre concrètement le RBAC dynamique~: déplacer une
responsabilité (accorder ou retirer une permission à un rôle) modifie simultanément le menu, les
routes accessibles et les boutons visibles, \textbf{sans aucune modification de code}.

\begin{longtable}{|m{3.4cm}|m{4.2cm}|m{6.2cm}|}
\hline
\textbf{Écran} & \textbf{Permission requise} & \textbf{Visibilité par défaut} \\
\hline
\endhead
\endfoot
\endlastfoot
\hline
Tableau de bord & (authentifié) & Tous \\
\hline
Projets & \texttt{VIEW\_PROJECT} & Directeur, Chef de projet, Développeur \\
\hline
Charges & \texttt{VIEW\_WORKLOAD} & Directeur, Chef de projet, Développeur \\
\hline
KPI & \texttt{VIEW\_KPI} & Directeur, Chef de projet \\
\hline
Facturation & \texttt{VIEW\_BILLING} & Directeur, Chef de projet \\
\hline
Missions & \texttt{VIEW\_MISSION} & Directeur, Chef de projet, Développeur \\
\hline
Gouvernance & \texttt{VIEW\_GOVERNANCE} & Directeur, Chef de projet \\
\hline
Ressources (TCC) & \texttt{VIEW\_RESOURCES} & Administrateur, Directeur, Chef de projet \\
\hline
Administration & \texttt{MANAGE\_USERS} & Administrateur \\
\hline
\caption{Correspondance écran $\rightarrow$ permission $\rightarrow$ visibilité (matrice corrigée)}
\label{tab:front-rbac}
\end{longtable}

\section{Communication avec l'API REST}

Chaque module dispose d'un service injectable typé qui encapsule les appels HTTP vers l'API,
dont l'adresse de base est centralisée dans le fichier d'environnement
(\texttt{environment.apiUrl}). Les échanges sont fortement typés par des interfaces TypeScript
miroir des DTO backend (\texttt{Project}, \texttt{Mission}, \texttt{JalonFacturation}\dots), ce qui
fait remonter à la compilation toute incohérence de contrat. L'état local des écrans est porté par
des \emph{Signals}, qui déclenchent un rafraîchissement ciblé de l'affichage à chaque mutation.

\section{Écrans métier}

\subsection{Cadre applicatif et tableau de bord}
Le composant \texttt{Shell} fournit l'ossature commune~: barre supérieure, menu latéral dynamique et
zone de contenu. Le tableau de bord présente une vue synthétique des projets accessibles à
l'utilisateur connecté.

\subsection{Projets et fiche d'identification}
La gestion de projet est l'écran le plus riche. Le formulaire de création/édition reproduit la
\textbf{fiche d'identification} du modèle Excel d'origine, organisée en trois sections~:

\begin{itemize}
  \item \textbf{Identification}~: code, nom, description, référence de contrat, client, bailleur de
        fonds, modèle business (seul / groupement), type d'engagement (forfait / régie) et, pour les
        utilisateurs habilités, le \textbf{chef de projet}.
  \item \textbf{Contractuel \& financier}~: statut, dates, devise et taux de conversion vers le
        dinar, budget initial, budget licences \& sous-traitance, provision pour pénalités.
  \item \textbf{Charge vendue}~: workload vendu et workload de garantie (en jours-homme).
\end{itemize}

\noindent Le formulaire affiche en temps réel les \textbf{grandeurs calculées} --- budget converti en
TND, provision pour risques (5~\%) et durée du contrat --- qui reproduisent exactement les valeurs de
la fiche Excel de référence (projet A24001~: 1\,269\,861~TND, PPR 63\,493~TND, durée 731~jours). Le
détail d'un projet est organisé en six onglets chargés à la demande~: informations \& KPI, équipe,
charges, facturation, missions et gouvernance.

\subsection{Saisie opérationnelle}
Les écrans de saisie s'appuient sur des fenêtres modales Bootstrap pilotées par des \emph{Signals}.
Ils couvrent l'ensemble du cycle métier~:

\begin{longtable}{|m{3.6cm}|m{10.4cm}|}
\hline
\textbf{Module} & \textbf{Actions d'écriture disponibles dans l'interface} \\
\hline
\endhead
\endfoot
\endlastfoot
\hline
Projets & Création/édition (fiche complète), assignation du chef de projet, gestion d'équipe (affecter / retirer) \\
\hline
Charges & Planification (plan de charge), saisie de charge réelle, validation \\
\hline
Facturation & Création de jalon, facturation d'un jalon, enregistrement de paiement, avenant (montant \emph{et} impact charge) \\
\hline
Missions & Création de mission et de ses composantes de coût (perdiem, billet, transport\dots) \\
\hline
Gouvernance & Risques, livrables (machine à états), demandes de changement, parties prenantes \\
\hline
Administration & Création / édition / désactivation d'utilisateurs \\
\hline
\caption{Couverture fonctionnelle des écrans de saisie}
\label{tab:front-forms}
\end{longtable}

\subsection{Ressources et tableaux de bord financiers}
L'écran \emph{Ressources} expose les tarifs journaliers (TCC) gérés par l'administrateur. Les
indicateurs financiers (budget planifié, consommé, EAC, marge, taux de consommation) sont présentés
sous forme de cartes et de barres de progression, réservés aux profils habilités (\texttt{VIEW\_KPI}).

\section{Reflet du périmètre de données}

L'interface respecte le périmètre d'autorisation appliqué par le backend (ADR-021)~: la liste des
projets et leurs ressources rattachées sont filtrées selon la relation de l'utilisateur au projet.
Le \textbf{directeur} (capacité d'accès portefeuille) voit l'ensemble des projets~; le \textbf{chef
de projet} ne voit que les projets qu'il pilote~; le \textbf{développeur} uniquement ceux auxquels il
est affecté. Toute tentative d'accès à un projet hors périmètre, même par saisie directe de l'URL,
est rejetée par le serveur (\texttt{403}), garantissant qu'aucune donnée d'un autre projet ne fuite
par l'interface.

\section*{Conclusion}
Le frontend Angular offre une interface complète, réactive et fidèle au modèle métier~: il reproduit
la fiche d'identification du classeur Excel, couvre l'intégralité du cycle de saisie opérationnelle
et applique le contrôle d'accès dynamique à trois niveaux (routes, menu, actions). Surtout, il
illustre la valeur centrale du projet~: parce que toute décision d'affichage dérive de l'ensemble des
permissions, une évolution de la politique d'autorisation se répercute instantanément sur l'interface
sans modification de code. Le chapitre suivant présente la stratégie de tests, la revue de sécurité et
la préparation au déploiement.
